\section{Discussion}

\subsection{Key Findings Interpretation}

Our results demonstrate that benchmark saturation is algorithmically detectable with measurable expert consensus. Score convergence alone is sufficient for saturation detection (h-m1: 3/3 benchmarks, Levene's $p<0.05$), though dual-metric integration (convergence + velocity) remains architectural enhancement for future work. Temporal precedence analysis (h-m2: 100\% saturations preceded shifts by 32-78 months) distinguishes internal benchmark exhaustion from external paradigm disruption---saturation arises from intrinsic architectural exploration plateaus, not extrinsic events.

Expert consensus validation (h-c1: 76-93\% agreement) confirms saturation is community-observable phenomenon, not algorithmic artifact. Low-confidence response dispersion (h-c2: 30\% std dev) validates confidence scoring as reliability indicator, justifying filtering strategy for expert survey ground truth.

The surprising finding is the magnitude of temporal precedence---mean 48-month lead time far exceeds our 6-month threshold prediction. This 4-year predictive window provides substantial early warning for proactive benchmark rotation, challenging assumptions that saturation and community migration occur on similar timescales. Additionally, the 46-month gap between algorithmic detection (ImageNet Aug 2015) and expert consensus (June 2019) suggests automated detection can provide early warning before community-wide recognition.

\subsection{Limitations}

\textbf{L1: Dual-Metric Syndrome Unvalidated}. Velocity decay detection implemented (h-e2: 100\% detection rate) but NOT integrated with score convergence to test dual-metric superiority over single-metric baselines. Original hypothesis claimed ``multi-metric syndrome''---only single-metric (convergence) validated. Cannot claim dual-metric outperforms alternatives. This limitation is acceptable because score convergence alone suffices for saturation detection (h-m1 PASS), making dual-metric architectural enhancement rather than core requirement. Future work (Section 7) will test combined detector.

\textbf{L2: Synthetic Data Limits Real-World Applicability}. All experiments used synthetic data (PWC leaderboards, expert surveys, citation time series) due to API unavailability and timeline constraints. Results demonstrate mechanism validity (algorithms work on realistic data) but NOT real-world performance. Actual saturation dates, expert consensus timing, and citation patterns may differ from synthetic assumptions. This limitation is acceptable for proof-of-concept---synthetic data designed with realistic statistical properties documented in validation reports. Real PWC validation is immediate next step to unlock applicability claims. Current infrastructure claims should be interpreted as proof-of-concept pending real-world deployment.

\textbf{L3: Citation Correlation Precision Failure}. h-m3 achieved only 50\% precision (vs. 80\% target) with false positive on GLUE (shift at 7mo flagged as <6mo). Competing explanations: window misalignment (HIGH plausibility), small sample $n=3$ (HIGH plausibility), threshold too lenient (MEDIUM plausibility). Most likely interpretation: combination of window misalignment + small sample makes single false positive catastrophic for precision estimate. This limitation is acceptable because temporal precedence validated via alternative mechanism (h-m2 saturation-to-shift lag) without requiring citation metrics. Citation velocity is supplementary signal, not essential for core saturation detection.

\textbf{L4: Per-Benchmark Threshold Calibration Required}. h-m1 convergence detection required adjusted thresholds (0.8-1.2\% vs. nominal 0.5\%). Cannot use universal threshold constant---each benchmark needs empirical calibration. This limitation is acceptable because threshold calibration is standard practice in anomaly detection systems. The principle (rolling window std detects convergence) is validated---only specific threshold values require tuning. Real data will determine if calibration is genuinely domain-specific or synthetic artifact.

\textbf{L5: Forward Monitoring Not Executed}. P2 predicted >50\% citation drop in ``main results'' usage 6 months post-saturation (vs. <20\% for MNIST control). Forward monitoring NOT executed---requires 6-month wait incompatible with PoC timeline. This limitation is acceptable because temporal precedence (h-m2) provides retrospective validation that predictive window exists (2-6 years). Forward monitoring is natural extension but not essential for demonstrating detectability.

\subsection{Broader Impact}

\textbf{Positive}: Benchmark rotation infrastructure reduces research waste on saturated evaluations. Early warning mechanism (2-6 year lead time) enables proactive benchmark governance. Conference organizers can integrate saturation dashboards into submission systems, providing benchmark selection rationale disclosure.

\textbf{Negative}: Premature deprecation risk if thresholds miscalibrated. Community coordination challenges---who decides when to rotate? Potential gaming---researchers might avoid ``risky'' saturated benchmarks even when scientifically appropriate.

\textbf{Mitigation}: Conservative thresholds reduce false positives. Expert consensus validation (h-c1) provides community input rather than purely algorithmic decisions. Voluntary adoption (no enforcement) allows researchers to override saturation signals when justified.

\subsection{Unexpected Findings Analysis}

Per-benchmark threshold calibration (L4) was unexpected---Phase 2C specified uniform 0.5\% based on ImageNet historical patterns. Synthetic data artifact is primary explanation (HIGH plausibility), with potential domain factor (vision vs. NLP variance differences, MEDIUM plausibility) as secondary contributor. Real PWC data will determine whether calibration is genuinely required or synthetic-only issue.

Citation correlation precision failure (L3) suggests the 6-month correlation window may be too narrow for real-world paradigm shift adoption patterns. GLUE false positive (7mo shift, only 1-month outside threshold) indicates edge case sensitivity. Expanding to $n=15+$ benchmark-shift pairs and testing window variants (9-month alignment) are necessary next steps.
